% Propietario preservado y actualizado por los cierres sucesores.
\section{Auto-inertie comme connexion endogène de l'état conjoint}

Soit \(\nabla^{\rm tot}\) une connexion sur l'espace des états
conjoints. Le cadre local de l'observateur s'obtient par restriction :
\begin{equation}
 \mathfrak F_{\mathcal O}(\Gamma)
 =\operatorname{Res}_{\mathcal O}
 (\Gamma,\nabla^{\rm tot}).
 \label{eq:rm-obs-induced-frame}
\end{equation}
Une trajectoire complète \(\Gamma(t)\) est auto-inertielle lorsque
\begin{equation}
 \nabla^{\rm tot}_{\dot\Gamma}\dot\Gamma=0.
 \label{eq:rm-obs-total-geodesic}
\end{equation}
Soit \(\pi_S\) la projection sur le sous-système observé et
\(\gamma_S=\pi_S\circ\Gamma\). Nous définissons le défaut de projection
\begin{equation}
 \mathcal K_{\pi_S}(\dot\Gamma,\dot\Gamma)
 :=\nabla^S_{\dot\gamma_S}\dot\gamma_S
 -\dd\pi_S\!\left(
 \nabla^{\rm tot}_{\dot\Gamma}\dot\Gamma
 \right).
 \label{eq:rm-obs-projection-defect}
\end{equation}
Sur une trajectoire auto-inertielle, il reste
\begin{equation}
 \boxed{
 a_S=\nabla^S_{\dot\gamma_S}\dot\gamma_S
 =\mathcal K_{\pi_S}(\dot\Gamma,\dot\Gamma).}
 \label{eq:rm-obs-projected-acceleration}
\end{equation}
L'accélération du sous-système mesure le couplage et les degrés de liberté
omis par la projection. La totalité possède une connexion endogène ;
ses ombres peuvent présenter une accélération, une rotation et une force effective.

\begin{definition}[État auto-inertiel]
Un état conjoint est auto-inertiel par rapport à une connexion et à une famille
d'observateurs lorsque son évolution satisfait
\cref{eq:rm-obs-total-geodesic} et que chaque cadre local s'obtient par
\cref{eq:rm-obs-induced-frame}. Les forces apparentes de chaque sous-système sont
enregistrées par les défauts \(\mathcal K_{\pi_S}\).
\end{definition}

La définition intègre l'observateur et l'observable dans une seule dynamique.
Le changement de cadre correspond à une autre restriction du même état et de la
même connexion. La distinction entre transport affine, rotation de la
tétrade, torsion et courbure reste disponible dans le codomaine
géométrique.

\section{Critère machien de factorisation globale}

Soit \(b:\mathcal X_\infty\to B_\infty\) la trace globale de frontière et soit
\(I_v:\mathcal X_\infty\to V_v\) une réponse inertielle locale. La
réponse est machienne lorsqu'il existe
\begin{equation}
 \overline I_v:\operatorname{im}b\longrightarrow V_v,
 \qquad
 I_v=\overline I_v\circ b.
 \label{eq:rm-obs-mach-factorization}
\end{equation}

\begin{theorem}[Critère de réponse machienne]
\label{thm:rm-obs-mach-criterion}
Il existe une unique application \(\overline I_v\) qui satisfait
\cref{eq:rm-obs-mach-factorization} si et seulement si \(I_v\) est constante
sur chaque fibre de \(b\) :
\begin{equation}
 b(x)=b(x')\quad\Longrightarrow\quad I_v(x)=I_v(x').
 \label{eq:rm-obs-mach-fiber-criterion}
\end{equation}
Si \(b\) est surjective, \(\overline I_v\) est définie sur tout
\(B_\infty\).
\end{theorem}

\begin{proof}
La factorisation implique la constance sur les fibres. Réciproquement, pour
\(\beta\in\operatorname{im}b\), nous définissons
\(\overline I_v(\beta)=I_v(x)\) avec \(x\in b^{-1}(\beta)\). La constance
garantit l'indépendance du représentant et l'équation de factorisation
garantit l'unicité.
\end{proof}

Le critère traduit le principe de Mach en une propriété vérifiable : la
réponse locale est une lecture de la classe globale de clôture. Pour démontrer
une dépendance véritablement globale, on exhibe en outre deux états de même
restriction locale et de traces de frontière distinctes dont les réponses diffèrent.
Cette condition sépare une factorisation globale effective d'une réponse
déterminée exclusivement par le voisinage local.

L'échelle gravitationnelle \(Q\) de l'horloge CT108 fournit une réalisation
adimensionnelle de cette lecture :
\begin{equation}
 \alpha_G(E_1,E_2)=\frac{E_1E_2}{Q^2}.
 \label{eq:rm-obs-global-clock-coupling}
\end{equation}
L'horloge globale fixe \(Q\), chaque section locale publie \(E/Q\), et la
géométrie répond par la compliance \(G/c^4\). L'équation
fournit une carte machienne de la réponse gravitationnelle ; le
\cref{thm:rm-obs-mach-criterion} détermine l'exigence de factorisation
que doit satisfaire sa réalisation sur l'état complet.

\section{Holonomie et hiérarchie des retours}

Soit \(\gamma\) un lacet basé en \(x_0\), et soit
\begin{equation}
 \mathscr H_\gamma=\Hol_\nabla(\gamma)
 \label{eq:rm-obs-holonomy}
\end{equation}
l'holonomie dans la fibre \(E_{x_0}\). Une donnée transportée revient en tant que
vecteur exactement lorsqu'elle appartient à
\begin{equation}
 \operatorname{Fix}(\mathscr H_\gamma)
 =\{\psi:\mathscr H_\gamma\psi=\psi\}.
 \label{eq:rm-obs-fix-holonomy}
\end{equation}
L'interprétation auto-inertielle se condense dans cette égalité : le cadre
transporté par la connexion de la totalité conserve la donnée située dans
le sous-espace fixe.

Pour une sous-algèbre d'observables accessibles
\(\mathcal A_{\mathcal O}\), nous définissons
\begin{align}
 \mathcal R_{\rm vec}
 &=\{\psi:\mathscr H_\gamma\psi=\psi\},
 \label{eq:rm-obs-return-vector}\\
 \mathcal R_{\rm ray}
 &=\{\psi:\mathscr H_\gamma\psi=\ee^{i\theta}\psi
 \text{ pour un certain }\theta\},
 \label{eq:rm-obs-return-ray}\\
 \mathcal R_{\rm obs}
 &=\{\psi:
 \langle\mathscr H_\gamma\psi,A\mathscr H_\gamma\psi\rangle
 =\langle\psi,A\psi\rangle
 \text{ pour tout }A\in\mathcal A_{\mathcal O}\}.
 \label{eq:rm-obs-return-observable}
\end{align}

\begin{proposition}[Hiérarchie vecteur--rayon--observable]
\label{prop:rm-obs-return-hierarchy}
On a
\begin{equation}
 \boxed{
 \mathcal R_{\rm vec}
 \subseteq\mathcal R_{\rm ray}
 \subseteq\mathcal R_{\rm obs}.}
 \label{eq:rm-obs-return-hierarchy}
\end{equation}
\end{proposition}

\begin{proof}
Un vecteur fixe est un rayon fixe de phase nulle. Si
\(\mathscr H_\gamma\psi=\ee^{i\theta}\psi\), les phases conjuguées s'
annulent dans toute espérance quadratique, de sorte que tous les observables
accessibles conservent leur valeur.
\end{proof}

La monodromie nonadique occupe naturellement cette hiérarchie :
\begin{equation}
 g(k+9)=g(k),
 \qquad
 B_{k+9}\neq B_k\quad\hbox{en général}.
 \label{eq:rm-obs-phase-state-return}
\end{equation}
Le retour de phase appartient à la lecture observable ; le changement de
\(B_k\) conserve l'incrément de mémoire. Une réalisation injective de l'
état complet dans la fibre permet d'élever le critère
\(\operatorname{Fix}(\mathscr H_\gamma)\) au retour généalogique total.
En l'absence de cette injectivité, le critère affirme le retour de la donnée
réalisée.

\section{Ruban de Möbius et conjugaison des orientations}

Sur un mot dodécaphasique
\(\mathbf t=(t_0,\ldots,t_{11})\), l'involution généalogique est
\begin{equation}
 C_M(\mathbf t)=-\operatorname{rev}(\mathbf t),
 \qquad
 C_M^2=I.
 \label{eq:rm-obs-word-involution}
\end{equation}
Soit \(\ell_{\rm ret}\) un lacet de retour enrichi dont le relèvement
possède le caractère de signe
\begin{equation}
 \varepsilon_x:\langle\ell_{\rm ret}\rangle\simeq\mathbb Z
 \longrightarrow\{\pm1\},
 \qquad
 \varepsilon_x(\ell_{\rm ret})=-1.
 \label{eq:rm-obs-sign-character}
\end{equation}
Le fibré en intervalles associé est
\begin{equation}
 \mathcal M_x
 =([0,1]\times[-1,1])/(0,u)\sim(1,-u).
 \label{eq:rm-obs-mobius-bundle}
\end{equation}

\begin{theorem}[Retour relatif de Möbius]
\label{thm:rm-obs-mobius-return}
Sous le relèvement compatible de \(C_M\), le lacet et le caractère de
\cref{eq:rm-obs-sign-character}, \(\mathcal M_x\) possède une première classe de
Stiefel--Whitney non nulle. Un tour échange les orientations locales et
deux tours rétablissent l'orientation de la section.
\end{theorem}

\begin{proof}
La fonction de transition du quotient
\cref{eq:rm-obs-mobius-bundle} est la représentation signe de
\(\pi_1(S^1)\). Sa classe \(w_1\) est le générateur non trivial de
\(H^1(S^1;\mathbb Z/2\mathbb Z)\). La première itération multiplie la
coordonnée de fibre par \(-1\) et la seconde par \((-1)^2=1\).
\end{proof}

La formulation conserve trois retours distincts : la clôture de phase
\(C_9\), le retour d'orientation dans le revêtement double et l'incrément de
mémoire dans le relèvement \(\mathbb Z\). Un tour peut rétablir le
rayon et changer le vecteur ; deux tours rétablissent l'orientation et peuvent
conserver une profondeur différente. Le ruban exprime précisément la
différence entre publication et généalogie.

Les secteurs pair et impair d'une fonction \(f\) sur le revêtement double sont
\begin{equation}
 f^+=\frac12(f+C_M^*f),
 \qquad
 f^-=\frac12(f-C_M^*f).
 \label{eq:rm-obs-even-odd-decomposition}
\end{equation}
La partie paire descend à la base ; la partie impaire est une section de la
droite de signe. La mémoire quadratique réside dans le secteur pair, tandis que l'
orientation et la charge linéaire résident dans le secteur impair.

\section{Conjugaison antiunitaire et orientation temporelle}

L'involution de mot \(C_M\), l'inversion centrale de l'atlas et la
conjugaison antiunitaire de la réalisation physique possèdent des domaines
distincts. Soit \(C\) la conjugaison antiunitaire sur une réalisation de
Hilbert, soit \(H=H^*\) l'hamiltonien et soit \(J_E\) la structure complexe
dans la réalification. Supposons
\begin{equation}
 CJ_EC^{-1}=-J_E,
 \qquad
 CHC^{-1}=H,
 \qquad
 J_EH\subset HJ_E,
 \label{eq:rm-obs-antiunitary-hypotheses}
\end{equation}
avec \(J_E^*=-J_E\), \(J_E^2=-I\) et préservation du domaine de \(H\).
La commutation exprime la linéarité complexe de la dynamique par rapport à
la structure \(J_E\).

\begin{proposition}[Renversement antiunitaire de l'évolution]
\label{prop:rm-obs-antiunitary-time}
Pour
\begin{equation}
 U(t)=\exp(-J_EHt/\hbar_{\rm int}),
 \label{eq:rm-obs-unitary-evolution}
\end{equation}
on a
\begin{equation}
 \boxed{CU(t)C^{-1}=U(-t).}
 \label{eq:rm-obs-time-reversal}
\end{equation}
\end{proposition}

\begin{proof}
Le générateur \(-J_EH/\hbar_{\rm int}\) est antiadjoint. Les hypothèses de
\cref{eq:rm-obs-antiunitary-hypotheses} changent son signe sous conjugaison.
Le calcul fonctionnel et l'unicité du groupe unitaire engendré produisent
\cref{eq:rm-obs-time-reversal}.
\end{proof}

Lorsque, en outre, un opérateur de charge \(Q_{\rm ch}\) et un observable de
masse \(M\) satisfont
\begin{equation}
 CQ_{\rm ch}C^{-1}=-Q_{\rm ch},
 \qquad
 CMC^{-1}=M,
 \label{eq:rm-obs-charge-mass-parity}
\end{equation}
les deux orientations du ruban sont publiées comme porteuses de charge
conjuguée et de mémoire paire partagée. La représentation particule--
antiparticule requiert l'entrelaceur qui relève \(C_M\) en \(C\) et
conserve l'évolution. Les identités
\cref{eq:rm-obs-word-involution,eq:rm-obs-time-reversal} fournissent les
deux extrémités typées de cet entrelaceur.

Le lecteur temporel tritique organise les régimes par
\begin{equation}
 \begin{aligned}
 +1&\longmapsto\text{retour, mémoire et passé},\\
 0&\longmapsto\text{seuil torsionnel et présent},\\
 -1&\longmapsto\text{ouverture bidirectionnelle et futur}.
 \end{aligned}
 \label{eq:rm-obs-temporal-trit}
\end{equation}
L'assignation est un lecteur de la structure algébrique
\(J^2=-1,0,+1\). La promotion en courbures d'un espace-temps exige une
connexion et une carte métrique ; le lecteur temporel conserve entre-temps
son sens exact de retour, de seuil et d'ouverture.

\section{Récupérabilité et principe de Landauer}

Le principe thermodynamique de Landauer localise le coût physique de l'
effacement logiquement irréversible \cite{Landauer1961}.

Soit \((\Omega,2^\Omega,\mu)\) un espace probabilisé fini, soit
\(X:\Omega\to\mathcal X\) la variable d'état complet et soit
\(q:\mathcal X\to Y\) un lecteur. La règle de chaîne fournit
\begin{equation}
 \mathsf H(X)
 =\mathsf H(q(X))+\mathsf H(X\mid q(X)).
 \label{eq:rm-obs-entropy-chain}
\end{equation}
Le second terme est l'information contenue dans la fibre du lecteur.

\begin{theorem}[Landauer nul du transport complet]
\label{thm:rm-obs-zero-landauer}
Si \(U:\mathcal X\to\mathcal X\) est une mise à jour bijective, alors
\begin{equation}
 \boxed{\mathsf H(X\mid U(X))=0.}
 \label{eq:rm-obs-zero-conditional-entropy}
\end{equation}
L'information logique effacée par la mise à jour est nulle. Pour tout
protocole physique \(\mathcal P_U\) qui la réalise, la borne de Landauer et sa
limite réversible idéale satisfont
\begin{equation}
 \boxed{
 Q_{\rm L}[\mathcal P_U]\geq0,
 \qquad
 \inf_{\mathcal P_U^{\rm rev}}Q_{\rm L}[\mathcal P_U]=0.}
 \label{eq:rm-obs-zero-landauer-cost}
\end{equation}
\end{theorem}

\begin{proof}
L'inverse de \(U\) reconstruit \(X\) à partir de \(U(X)\) ; la
distribution conditionnelle est donc concentrée sur un seul état. Son entropie
est nulle et l'opération logique conserve l'information. La borne
thermodynamique obtenue est \(Q_{\rm L}\geq0\) ; les protocoles
quasi statiques réversibles réalisent l'infimum idéal nul. Cette affirmation
distingue la borne de la dissipation concrète d'un dispositif.
\end{proof}

La sortie projetée \(q(U(X))\) conserve une partie de l'information et
laisse dans la fibre
\begin{equation}
 \mathsf H(X\mid q(U(X))).
 \label{eq:rm-obs-reader-forgotten-info}
\end{equation}
Soit \(M:\mathcal X\to\mathfrak M\) un registre de mémoire et définissons le
lecteur enrichi
\begin{equation}
 \widehat q(x)=(q(x),M(x)).
 \label{eq:rm-obs-enriched-reader}
\end{equation}

\begin{proposition}[Récupérabilité par la sortie et la mémoire]
\label{prop:rm-obs-recoverability}
Si \(\widehat q\) est injectif sur le support de \(X\), alors
\begin{equation}
 \boxed{\mathsf H(X\mid q(X),M(X))=0.}
 \label{eq:rm-obs-recoverability}
\end{equation}
La quantité \(\mathsf H(X\mid q(X))\) mesure la mémoire supplémentaire requise
pour retrouver l'état à partir de la sortie visible.
\end{proposition}

\begin{proof}
L'injectivité sur le support fournit un inverse déterministe de
\(\widehat q(X)\) vers \(X\). L'entropie conditionnelle est donc nulle. La
règle de chaîne localise la différence entre le lecteur visible et le
lecteur enrichi dans la fibre de \(q\).
\end{proof}

Pour la réinitialisation d'un trit uniforme, l'entropie logique est \(\log3\),
et la borne de Landauer prend la forme
\begin{equation}
 \boxed{Q_{\rm reset}\geq k_B\Theta\log3.}
 \label{eq:rm-obs-trit-reset}
\end{equation}
L'identité structurelle de l'horloge,
\begin{equation}
 E_P=k_B\Theta_{\rm clk}\log3,
 \label{eq:rm-obs-planck-trit-landauer}
\end{equation}
évalue cette borne à la température chronologique déclarée. Le quantum de
Planck apparaît alors comme le coût minimal de réinitialisation d'une décision
ternaire uniforme dans cette carte thermique.

Le Landauer nul et le coût \(k_B\Theta\log3\) décrivent des opérations
différentes. Le premier correspond au transport bijectif de l'état
enrichi ; le second correspond à une réinitialisation qui identifie trois
registres. Les coûts de contrôle, de vitesse finie, de préparation et de
couplage à l'environnement appartiennent à la réalisation physique du
dispositif. L'égalité logique de coût minimal nul conserve l'
information ; une équation du mouvement pour un rebond ou un cycle
cosmologique requiert en outre une dynamique de frontière et ses bilans.

\section{Mesure réversible, conditionnement et réinitialisation}

Les opérations du processus complet s'ordonnent en
\begin{equation}
 \boxed{\begin{gathered}
 \text{préparation}
 \longrightarrow
 \text{couplage réversible}
 \longrightarrow
 \text{lecture du pointeur}\\
 \longrightarrow
 \text{conditionnement de la fibre}
 \longrightarrow
 \text{réinitialisation}
 \end{gathered}}
 \label{eq:rm-obs-measurement-chain}
\end{equation}
Chaque flèche possède son propre bilan informationnel :
\begin{enumerate}[label=\textup{(\roman*)}]
 \item La prémesure unitaire conserve la pureté et la norme dans le système
 composé.
 \item La lecture publie \(q(X)\) et conserve le registre \(M(X)\) dans l'
 appareil.
 \item Le conditionnement sélectionne la fibre relative et normalise sa
 mesure.
 \item Le canal non sélectif somme tous les instruments et préserve la
 trace.
 \item La réinitialisation identifie des registres et active la borne de
 Landauer correspondant à l'information effacée.
\end{enumerate}

Cette chaîne localise l'origine de l'irréversibilité opérative. L'état
global peut évoluer isométriquement tandis qu'une lecture partielle
présente une entropie croissante. L'information accessible à une branche et l'
information conservée dans les corrélations globales sont des grandeurs
distinctes. L'observateur couplé fait partie de ces corrélations et sa
mémoire constitue une composante physique de l'état conjoint.

\section{Confluence entre observation, auto-inertie et réponse modulaire}

Les résultats précédents s'assemblent dans la chaîne causale
\begin{equation}
 \begin{aligned}
 \text{état APP--TRIT--TPK enrichi}
 &\longrightarrow
 (\text{section},\text{observateur},\text{mémoire})\\
 &\longrightarrow
 (\text{fibre},\text{PVM},\text{instrument})\\
 &\longrightarrow
 (\text{état relatif},\text{holonomie},\text{repère induit})\\
 &\longrightarrow
 (\text{réponse machienne},\text{Landauer typé}).
 \end{aligned}
 \label{eq:rm-obs-causal-chain}
\end{equation}
La structure modulaire du chapitre précédent ajoute la paire
\begin{equation}
 (\widehat{\mathcal A},K_A),
 \label{eq:rm-obs-area-modular-pair}
\end{equation}
qui conserve la quantité géométrique et la provenance. L'observateur accède à une
compression de cette paire par sa PVM ; la mémoire conserve le canal dont
provient le résultat ; l'auto-inertie transporte le cadre le long de
la connexion ; l'entropie relative mesure le coût du changement d'état
généalogique ; et Landauer localise le coût de l'effacement effectif du
registre.

La relation avec Everett se situe au premier niveau de cette chaîne : l'
isométrie conserve la superposition des registres et l'observateur occupe un
état relatif. La relation avec Mach se situe au niveau de la connexion :
le cadre et la réponse locale sont induits par la totalité compatible. Le
ruban de Möbius se situe au niveau de l'holonomie : l'orientation
peut changer tandis que le rayon et les observables accessibles reviennent. Les
trois lectures partagent l'état enrichi et conservent des fonctions
mathématiques différentes.

\section{Domaine de validité et conditions de réalisation physique}

La force probatoire se répartit comme suit :
\begin{enumerate}[label=\textup{(\roman*)}]
 \item La correspondance fibre--PVM--Lüders, la représentation de
 Kraus, le canal CPTP et la dilatation sont des
 \status{résultats retrouvés exacts} dans les domaines déclarés.
 \item La formulation de l'état relatif et la conservation isométrique
 du canal sont \status{exactes}. L'ontologie everettienne des branches et leur
 stabilité macroscopique relèvent d'une
 \status{interface interprétative et dynamique}.
 \item Le critère
 \(\operatorname{Fix}(\Hol)\), la hiérarchie
 vecteur--rayon--observable et la factorisation machienne sont
 \status{résultats retrouvés exacts}. Leur application à une géométrie
 gravitationnelle exige le morphisme de réalisation correspondant.
 \item Le ruban de Möbius est \status{exacta relativamente} au
 relèvement, au lacet et au caractère de signe déclarés. L'
 interprétation particule--antiparticule exige l'entrelaceur
 antiunitaire de \cref{eq:rm-obs-antiunitary-hypotheses}.
 \item Le Landauer nul d'une mise à jour bijective, la récupérabilité
 du lecteur enrichi et la borne de la réinitialisation ternaire sont
 \status{résultats retrouvés exacts}. Le sélecteur cosmologique et la
 dynamique des feuillets prolongent ces bilans par les équations d'
 évolution, les conditions de raccordement et la mémoire torsionnelle déclarées.
\end{enumerate}

\begin{proofstatusbox}
Le chapitre démontre dans ses domaines finis l'équivalence des
conditionnements, la conservation du canal global, les critères de
retour et de factorisation, et les bilans logiques de Landauer. La
réalisation conjointe emploie l'entrelaceur isométrique entre la mémoire TPK et
l'appareil, la stabilité dynamique du registre fini et le transport de l'
involution généalogique vers la conjugaison antiunitaire avec ses observables pairs
et impairs.
\end{proofstatusbox}

\begin{falsifierbox}
Réfutent les énoncés correspondants : une fibre non mesurable ; une famille
de projecteurs dépourvue de PVM commune ; un conditionnement de Lüders
qui diffère de \cref{eq:rm-obs-luders-fiber-equivalence} ; des opérateurs de
Kraus dont la somme des effets diffère de l'identité ; une dilatation avec
\(V_a^\dagger V_a\neq I\) ; un retour vectoriel exact attribué à une donnée
hors de \(\operatorname{Fix}(\Hol)\) ; une réponse machienne variable au sein d'une
fibre de \(b\) ; un lacet de caractère de signe trivial présenté comme un ruban
de Möbius ; une conjugaison qui ne satisfait pas
\(CU(t)C^{-1}=U(-t)\) ; une mise à jour bijective avec une entropie
conditionnelle positive ; ou une réinitialisation ternaire uniforme dont le bilan omet l'
information \(\log3\). Réfute également une promotion cosmologique la
présentation du Landauer nul comme équation dynamique de rebond.
\end{falsifierbox}
